% Contenido del póster: Danisewicz y Elard (2023), Journal of Banking & Finance. DOI 10.1016/j.jbankfin.2023.106986
% Elaborado con apoyo de IA a partir del resumen oficial del artículo (OpenAlex), 2026-11-12.
\def\DOI{10.1016/j.jbankfin.2023.106986}
\def\TituloEs{Los efectos reales de la tecnología financiera: préstamos por plataformas y quiebra personal}
\def\TituloOriginal{The real effects of financial technology: Marketplace lending and personal bankruptcy}
\def\Autores{Danisewicz y Elard}
\def\Anio{2023}
\def\Revista{Journal of Banking \& Finance}

\def\Problema{Se discute si el crédito de las plataformas digitales ayuda o perjudica a los hogares. Medir su efecto es difícil porque el acceso a ese crédito no es aleatorio.}
\def\Objetivo{Examinar cómo la tecnología financiera afecta las dificultades económicas de los hogares, medidas por las quiebras personales.}
\def\Metodo{Los autores aprovechan un fallo judicial como fuente externa de variación: los préstamos con tasas por encima del límite de usura otorgados por bancos a residentes de Connecticut y Nueva York quedaban anulados si se vendían a entidades no bancarias. Comparan la evolución de las quiebras personales y del crédito por plataformas antes y después del fallo.}
\def\Resultados{\begin{itemize}
  \item Tras el fallo, las quiebras personales aumentaron de forma persistente.
  \item El crédito por plataformas disminuyó, sobre todo en los hogares de bajos ingresos.
  \item Los impagos de los préstamos de plataforma y el resto del crédito al consumo de bancos y financieras no cambiaron.
\end{itemize}}
\def\Conclusion{El aumento de las quiebras se explica principalmente por la pérdida de acceso a una nueva tecnología crediticia: las plataformas cumplían un papel en el financiamiento de los hogares de bajos ingresos.}
\def\Relevancia{En el Perú, los proyectos personales fueron una parte pequeña del financiamiento participativo: se publicaron \cifra{publicadosPer} y casi todos correspondieron a una plataforma que dejó de registrar financiamiento en \cifra{finTercera}; tras la modificación de 2025 no se publicó ninguno. El estudio sugiere que la desaparición de ese segmento puede tener costos para los hogares.}
\def\PalabrasClave{préstamos por plataformas, quiebra personal, usura, hogares de bajos ingresos, tecnología financiera}
% --- Ficha del estudio (herramientas/posteres/ficha_articulos.py) ---
\def\Ambito{Estados Unidos (Connecticut y Nueva York)}
\def\Enfoque{Experimento natural: fallo judicial sobre préstamos por encima de la usura}
\def\FiguraPeru{2}
% --- fin de la ficha ---
